% This file was converted to LaTeX by Writer2LaTeX ver. 1.9.9
% see http://writer2latex.sourceforge.net for more info
\documentclass{article}
\usepackage{calc,amsmath,amssymb,amsfonts}
\usepackage[LGR,T1]{fontenc}
\usepackage[greek,italian]{babel}
\usepackage[style=numeric,backend=biber]{biblatex}
\author{utente}
\date{2026-08-18}
\begin{document}
Capitolo 1 – Il computer non si stanca

La settimana che separava una lezione dalla successiva mi sembrò molto più lunga di sette giorni. A quell'età il tempo
aveva ancora la capacità di dilatarsi attorno alle cose che aspettavo: una domenica, una festa, un regalo promesso
potevano sembrare lontanissimi, e quella volta ad attendermi non c'era nemmeno qualcosa che avrei potuto prendere in
mano una volta tornato a casa. Non possedevo ancora un computer sul quale ripetere gli esperimenti della prima lezione;
avevo soltanto il ricordo di ciò che avevo visto e poche parole annotate nella memoria con una precisione sorprendente:
READY, PRINT, RUN.

Continuavo a ripensare soprattutto a quel CIAO apparso sul monitor. Era difficile spiegare perché cinque lettere mi
avessero impressionato tanto, considerando che avevo già visto videogiochi infinitamente più spettacolari. La
differenza, però, cominciava a diventarmi chiara: nei giochi potevo scegliere che cosa fare all'interno di regole
stabilite da qualcun altro, mentre in quell'aula avevo visto nascere una regola davanti ai miei occhi. Prima il
Commodore non scriveva nulla; poi qualcuno aveva digitato un'istruzione e la macchina aveva obbedito. Da lì le domande
avevano cominciato a moltiplicarsi: se potevamo fargli scrivere una parola, avremmo potuto farla comparire in un punto
preciso? Cambiarne il colore? Ripeterla? Far muovere qualcosa? E, inevitabilmente, tornava anche il pensiero ai
videogiochi e a quella vecchia domanda che non aveva mai smesso di accompagnarmi.

Quando arrivò finalmente il pomeriggio della seconda lezione, entrai nella scuola con una sensazione diversa rispetto
alla prima volta. Non ero più il bambino che varcava una soglia sconosciuta chiedendosi che cosa avrebbe trovato.
Conoscevo il percorso fino all'aula, sapevo dove Rossi Brunori avrebbe sistemato il Commodore e riuscivo quasi ad
anticipare il momento in cui il monitor si sarebbe acceso. Era bastato un solo incontro perché quella stanza, che una
settimana prima mi sembrava appartenere a un mondo più grande del mio, cominciasse già ad avere qualcosa di familiare.

Il professore era arrivato prima di noi e, insieme a suo figlio, stava terminando di collegare il computer, il monitor e
il drive. Guardavo quei gesti ripetersi quasi nello stesso ordine della volta precedente e mi accorgevo che stavano già
diventando una specie di rituale: i cavi sistemati sulla cattedra, il monitor orientato verso la classe, il Commodore
al centro e infine l'interruttore. Quando lo schermo azzurro comparve nuovamente e sotto le righe iniziali apparve
READY., ebbi la sensazione curiosa di ritrovare qualcosa che conoscevo da molto più di una settimana.

Rossi Brunori si voltò verso di noi e, con un sorriso, volle verificare quanto fosse sopravvissuto alla distanza tra le
due lezioni. $\text{\textgreek{«}}$Vediamo un po'. Qualcuno si ricorda come si chiamava il comando che abbiamo usato
per far scrivere una parola al Commodore?$\text{\textgreek{»}}$ Questa volta le mani si alzarono molto più velocemente
e PRINT arrivò quasi in coro. Alla domanda successiva, quella sul comando necessario per eseguire il programma, la
risposta fu altrettanto immediata: RUN.

$\text{\textgreek{«}}$Bene$\text{\textgreek{»}}$, commentò. $\text{\textgreek{«}}$Allora non avete imparato soltanto due
parole. Qualcosa avete capito.$\text{\textgreek{»}}$

Senza aggiungere altre spiegazioni, digitò il piccolo programma della settimana precedente:

10 PRINT {\textquotedbl}CIAO{\textquotedbl}

Premette RETURN, scrisse RUN e il saluto ricomparve sullo schermo, identico a quello che avevo continuato a ricordare
per tutta la settimana. Rossi Brunori lo osservò per un istante e poi disse una cosa che mi sorprese:
$\text{\textgreek{«}}$Adesso abbiamo un problema.$\text{\textgreek{»}}$

Il programma aveva appena fatto esattamente ciò che gli avevamo chiesto e non riuscivo a capire dove potesse nascondersi
il problema. Il professore indicò la parola CIAO e ci invitò a cambiare punto di vista. Il Commodore sapeva salutarci,
certo, ma lo faceva una sola volta. Se quella scritta fosse stata un messaggio da mostrare continuamente,
un'informazione destinata a ripetersi o semplicemente qualcosa che volevamo vedere comparire molte volte, avremmo
dovuto forse scrivere cento righe identiche?

L'idea fece sorridere qualcuno, ma era proprio l'assurdità di quella soluzione a rendere evidente il problema. Rossi
Brunori ci propose allora un esempio lontano dal computer: immaginate, disse, di stare all'ingresso di una scuola e di
dover salutare ogni studente che entra. Al primo direste $\text{\textgreek{«}}$ciao$\text{\textgreek{»}}$ senza
pensarci, al decimo comincereste probabilmente ad annoiarvi e al centesimo avreste già perso ogni entusiasmo. Una
macchina, invece, non avrebbe fatto alcuna differenza tra la prima e la millesima ripetizione. Non avrebbe provato
noia, non avrebbe perso concentrazione e non avrebbe chiesto di fare una pausa.

Quella spiegazione mi colpì perché spostava ancora una volta il confine tra ciò che consideravo una debolezza della
macchina e ciò che invece poteva diventare la sua forza. Il Commodore non capiva il significato di un saluto e non
provava alcun piacere nel ripeterlo, ma proprio per questo poteva farlo tutte le volte che gli avessimo chiesto. Il
problema non era convincerlo a lavorare di più; bisognava soltanto trovare il modo di dirgli che, dopo aver eseguito
un'istruzione, avrebbe dovuto ricominciare.

Rossi Brunori tornò davanti alla tastiera. $\text{\textgreek{«}}$Guardiamo il programma come lo guarda lui. Parte dalla
linea dieci, legge PRINT {\textquotedbl}CIAO{\textquotedbl}, esegue l'istruzione e poi cerca quella successiva. Che
cosa trova?$\text{\textgreek{»}}$

Dopo qualche esitazione qualcuno rispose: $\text{\textgreek{«}}$Niente.$\text{\textgreek{»}}$

$\text{\textgreek{«}}$Esatto. E allora si ferma. Non perché sia stanco, non perché ritenga che un saluto sia sufficiente
e nemmeno perché abbia deciso che il programma è finito. Si ferma perché noi non gli abbiamo dato altro da
fare.$\text{\textgreek{»}}$

Questa distinzione, così semplice da sembrare quasi ovvia una volta pronunciata, conteneva qualcosa che avrei incontrato
infinite volte negli anni successivi: il computer non completa i nostri pensieri. Arriva esattamente fino al punto in
cui arrivano le istruzioni e non un passo oltre. Se volevamo che il programma ricominciasse, dovevamo essere noi a
indicargli dove tornare.

Sul monitor comparve una seconda riga:

20 GOTO 10

Il professore non ne spiegò immediatamente il significato. Lasciò che la osservassimo e ci chiese di provare a
interpretarla. Qualcuno pronunciò GOTO come un'unica parola, finché Rossi Brunori non la separò lentamente in GO e TO:
vai a. A quel punto il numero dieci che seguiva smise di essere un dettaglio e diventò la destinazione dell'istruzione.
Alla linea venti stavamo dicendo al Commodore di tornare alla linea dieci.

Adesso il programma era formato da appena due righe:

10 PRINT {\textquotedbl}CIAO{\textquotedbl}

20 GOTO 10

Rossi Brunori seguì con un dito il percorso delle istruzioni: linea dieci, poi linea venti, di nuovo linea dieci, ancora
la venti. Continuò a ripetere quella sequenza con calma, quasi volesse costringerci a vedere il programma muoversi
prima ancora di eseguirlo, finché un ragazzo lo interruppe con una constatazione che sembrò arrivargli improvvisamente:
$\text{\textgreek{«}}$Ma allora non finisce più.$\text{\textgreek{»}}$

Il professore si fermò e lo guardò con un'espressione soddisfatta. $\text{\textgreek{«}}$Esatto.$\text{\textgreek{»}}$

Per qualche secondo non disse altro. Aveva ottenuto ciò che cercava: la conclusione non era arrivata da una definizione
scritta alla lavagna, ma dal tentativo di seguire mentalmente il comportamento del programma. Senza conoscerne ancora
il nome avevamo appena incontrato uno dei concetti fondamentali della programmazione, quello che oggi chiameremmo senza
esitazione un ciclo. Allora, per comprenderlo, ci erano bastate due righe e il percorso immaginario tracciato da un
dito davanti allo schermo.

Prima di eseguire il programma Rossi Brunori, però, introdusse un'altra domanda. Fino a quel momento avevamo cercato il
modo di convincere il computer a continuare; adesso che sembravamo esserci riusciti, bisognava domandarsi come avremmo
fatto a fermarlo. Era una caratteristica del suo modo di insegnare che avrei imparato a riconoscere: una soluzione
raramente chiudeva un argomento. Più spesso apriva la porta al problema successivo.

$\text{\textgreek{«}}$Vediamo prima se abbiamo ragionato bene$\text{\textgreek{»}}$, disse tornando alla tastiera.

Digitò RUN e premette RETURN.

Questa volta il risultato non ebbe nulla della quiete della prima lezione. La parola CIAO comparve una volta, poi una
seconda, poi ancora, e in pochi istanti cominciò a riempire il monitor con una velocità che nessuno di noi si
aspettava. Le righe scorrevano verso l'alto mentre il Commodore continuava a eseguire le due istruzioni senza
esitazione, trasformando quel semplicissimo saluto in una cascata di testo che sembrava non voler terminare.

Nell'aula si levò un brusio divertito. Qualcuno rise, altri si sporsero leggermente in avanti per osservare meglio lo
schermo. Io continuavo a guardare quelle parole moltiplicarsi e pensavo alla differenza enorme che poteva produrre una
sola riga aggiunta al programma. Una settimana prima mi era sembrato straordinario ottenere un CIAO; adesso, con un
GOTO, ne avevamo creati così tanti da non riuscire più a contarli.

Rossi Brunori lasciò che il programma continuasse ancora per qualche secondo, abbastanza a lungo da farci capire che il
Commodore non avrebbe preso spontaneamente alcuna decisione. Poi allungò una mano verso la tastiera e premette
RUN/STOP. Il flusso si interruppe immediatamente e, dopo tutto quel movimento, lo schermo immobile produsse quasi la
sensazione di un silenzio improvviso.

Il professore indicò il tasto appena utilizzato e ci spiegò che non era importante soltanto sapere come interrompere un
programma. Quell'esperimento mostrava qualcosa di più generale: se impartiamo a un computer un'istruzione che lo porta
a ripetere la stessa operazione, la macchina continuerà a farlo finché le istruzioni non prevederanno una via d'uscita
o qualcuno non interverrà dall'esterno. Non si fermerà perché pensa di aver lavorato abbastanza, non si accorgerà che
il risultato è diventato inutile e non proverà a interpretare ciò che avremmo preferito ottenere.

Poi guardò lo schermo ancora pieno di CIAO e sorrise. $\text{\textgreek{«}}$Per questo, quando un programma fa qualcosa
che non vi aspettavate, prima di prendervela con il computer conviene guardare quello che avete scritto. Molto spesso
da qualche parte c'è un errore di sbaglio.$\text{\textgreek{»}}$

Ormai quella strana espressione non ci sorprese come la prima volta. Stava già diventando parte del lessico del corso,
ma dietro il sorriso con cui la accoglievamo cominciava a formarsi un'idea molto meno scherzosa: un errore non era la
dimostrazione che il computer non funzionava e nemmeno qualcosa di cui vergognarsi. Era una differenza tra ciò che
avevamo immaginato e ciò che avevamo realmente scritto, e proprio quella differenza poteva insegnarci qualcosa.

Negli anni successivi avrei trascorso migliaia di ore davanti a programmi che facevano esattamente ciò che avevo scritto
e, proprio per questo, qualcosa di diverso da ciò che avevo intenzione di ottenere. Avrei imparato altri linguaggi,
lavorato con macchine infinitamente più potenti e incontrato errori molto più difficili da individuare di un GOTO che
tornava alla linea sbagliata. Eppure il principio non sarebbe cambiato: prima di sospettare della macchina, bisognava
tornare alle istruzioni.

Uscendo dall'aula quel pomeriggio non avrei saputo formulare una considerazione così precisa. Sapevo soltanto che due
righe di BASIC erano riuscite a trasformare un singolo saluto in qualcosa che sembrava non avere fine e che, per la
seconda volta in due lezioni, ero entrato con una domanda e ne uscivo con altre ancora. Invece di scoraggiarmi, quella
sensazione mi attirava sempre di più.

Il Commodore continuava a sembrarmi una macchina straordinaria, ma cominciavo lentamente a capire che la parte davvero
interessante non era ciò che sapeva fare da solo. Era ciò che potevamo imparare a chiedergli.


\bigskip
\end{document}
